\section{הרצאה: פיתוח משוואת הדיפוזיה ומשוואת הרצף}

בשיעור זה נעבור ממשוואות היפרבוליות (כמו משוואת הגלים) למשפחה חדשה של משוואות: משוואות פרבוליות, ובראשן \textbf{משוואת הדיפוזיה} (או משוואת החום). נפתח את המשוואה בשתי גישות שונות: גישה מאקרוסקופית המבוססת על חוקי שימור ורצף, וגישה מיקרוסקופית המבוססת על הלוך אקראי (\LR{Random Walk}).

\subsection{משוואת הרצף וחוק שימור האנרגיה}

כדי לפתח את משוואת הדיפוזיה, נתבונן בנפח שרירותי $V$ המלא בחומר (גז, נוזל או מוצק). נגדיר את האנרגיה הפנימית הכוללת בנפח זה כ-$U$.

\begin{definitionbox}{צפיפות אנרגיה פנימית}
נסמן את צפיפות האנרגיה הפנימית ב-$u(\vec{r},t)$. הקשר לאנרגיה הכוללת הוא:
\[
U(t) = \int_V u(\vec{r},t) \, dV
\]
\end{definitionbox}

לפי החוק הראשון של התרמודינמיקה, השינוי באנרגיה הפנימית בזמן, $\frac{dU}{dt}$, נובע משני גורמים:
\begin{enumerate}
    \item \textbf{עבודה ($W$):} כיווץ/התפשטות הנפח, עבודה חשמלית או כימית.
    \item \textbf{חום ($Q$):} מעבר אנרגיה ספונטני הנובע מהפרשי טמפרטורות (ללא עבודה יזומה).
\end{enumerate}

בפיתוח הנוכחי, נניח שהמערכת \textbf{אינה מבצעת עבודה} (נפח קבוע, ללא ריאקציות), ונתמקד בזרימת חום ספונטנית.

\subsubsection{פיתוח משוואת הרצף}
נגדיר את וקטור שטף החום $\vec{J}$. וקטור זה מתאר את כמות האנרגיה העוברת ליחידת שטח ליחידת זמן. כיוון הווקטור מצביע על כיוון הזרימה.

קצב השינוי של האנרגיה בתוך הנפח $V$ שווה לשטף האנרגיה \textbf{הנכנס} דרך השפה $\partial V$ (המעטפת):
\[
\frac{dU}{dt} = -\oint_{\partial V} \vec{J} \cdot \hat{n} \, dA
\]
הסימן השלילי נובע מכך ש-$\hat{n}$ הוא נורמל המצביע \textbf{החוצה} מהנפח. אם $\vec{J}$ מקביל ל-$\hat{n}$, אנרגיה יוצאת והנגזרת שלילית.

נשתמש בהגדרת $U$ ובמשפט הדיברגנץ (\LR{Gauss's Theorem}) כדי להמיר את אינטגרל המשטח לאינטגרל נפחי:
\[
\frac{d}{dt} \int_V u \, dV = - \int_V (\nabla \cdot \vec{J}) \, dV
\]
נניח שהחומר במנוחה מאקרוסקופית (אין הסעה, $\vec{v}=0$), ולכן ניתן להכניס את הנגזרת בזמן לתוך האינטגרל כנגזרת חלקית:
\[
\int_V \left( \frac{\partial u}{\partial t} + \nabla \cdot \vec{J} \right) \, dV = 0
\]
מכיוון שהשוויון נכון לכל נפח שרירותי $V$ (אנו יכולים לכווץ את הנפח לנקודה), האינטגרנד חייב להתאפס זהותית.

\begin{resultbox}{משוואת הרצף (Continuity Equation)}
\[
\boxed{\frac{\partial u}{\partial t} + \nabla \cdot \vec{J} = 0}
\]
משוואה זו היא חוק שימור מקומי: השינוי בצפיפות הגודל (אנרגיה/מסה/מטען) בנקודה שווה לדיברגנץ (ההתבדרות) של השטף באותה נקודה.
\end{resultbox}

\subsection{חוק פורייה וקבלת משוואת הדיפוזיה}

משוואת הרצף היא משוואה כללית מאוד. כדי לקבל את משוואת הדיפוזיה הספציפית לחום, עלינו להניח שתי הנחות פיזיקליות נוספות ("יחסי מבנה" - \LR{Constitutive Relations}):

\subsubsection{1. קשר בין אנרגיה לטמפרטורה}
נניח שהשינוי באנרגיה הפנימית נובע בעיקר משינוי בטמפרטורה (הנחה טובה לגזים ומוצקים, פחות לנוזלים דחוסים).
\[
u \approx c_v T
\]
כאשר $c_v$ הוא קיבול החום הסגולי (ליחידת נפח) ו-$T$ הטמפרטורה.
לכן: $\frac{\partial u}{\partial t} = c_v \frac{\partial T}{\partial t}$.

\subsubsection{2. חוק פורייה (Fourier's Law)}
אנו יודעים אינטואיטיבית שחום זורם מאזור חם לאזור קר. הקשר בין שטף החום לגרדיאנט הטמפרטורה הוא לינארי (עבור הפרשי טמפרטורה קטנים):

\begin{definitionbox}{חוק פורייה}
\[
\vec{J} = -\kappa \nabla T
\]
כאשר:
\begin{itemize}
    \item $\kappa$ (\LR{Kappa}): מקדם הולכת החום (\LR{Thermal Conductivity}). יחידות: $[\text{W}\cdot\text{m}^{-1}\cdot\text{K}^{-1}]$.
    \item הסימן המינוס: מבטיח שהזרימה היא במורד הגרדיאנט (מהחם לקר).
\end{itemize}
\end{definitionbox}

\subsubsection{קבלת משוואת הדיפוזיה}
נציב את שני הקשרים הנ"ל למשוואת הרצף:
\[
c_v \frac{\partial T}{\partial t} + \nabla \cdot (-\kappa \nabla T) = 0
\]
\[
c_v \frac{\partial T}{\partial t} = \nabla \cdot (\kappa \nabla T)
\]
בהנחה שהחומר הומוגני ($\kappa$ קבוע במרחב) וקיבול החום קבוע:
\[
\frac{\partial T}{\partial t} = \frac{\kappa}{c_v} \nabla^2 T
\]

\begin{resultbox}{משוואת הדיפוזיה (משוואת החום)}
נסמן את \textbf{מקדם הדיפוזיה} ב-$D = \frac{\kappa}{c_v}$ (יחידות: $[\text{m}^2/\text{s}]$):
\[
\boxed{\frac{\partial T}{\partial t} = D \nabla^2 T}
\]
\end{resultbox}

\begin{notebox}{משמעות מקדם הדיפוזיה}
היחידות של $D$ הן $[\text{Length}^2 / \text{Time}]$. זה מרמז שבדיפוזיה, המרחק האופייני שחומר/חום מתפשט גדל כמו \textbf{שורש הזמן} ($x \sim \sqrt{Dt}$), בשונה מתנועה גלית או בליסטית שבה $x \sim vt$. התהליך הוא איטי יותר.
\end{notebox}

\subsection{פיתוח מיקרוסקופית: הלוך אקראי (Random Walk)}

כעת נפתח את אותה משוואה מתוך שיקולים הסתברותיים, ללא הנחות תרמודינמיות, אלא מתוך מודל של "שיכור" ההולך על סריג.

\textbf{המודל:}
\begin{itemize}
    \item סריג חד-ממדי עם מרווחים $L$.
    \item בכל צעד זמן $\Delta t$, החלקיק זז ימינה בסיכוי $1/2$ או שמאלה בסיכוי $1/2$.
    \item נגדיר $P(x,t)$: ההסתברות (או הצפיפות) למצוא את החלקיק במיקום $x$ בזמן $t$.
\end{itemize}

\begin{center}
\begin{tikzpicture}
    % Draw lattice line
    \draw[<->] (-3.5,0) -- (3.5,0);
    
    % Draw lattice points
    \foreach \x in {-2, 0, 2}
        \draw[fill=black] (\x,0) circle (2pt);
    
    % Labels
    \node[below] at (0,-0.2) {$x$};
    \node[below] at (-2,-0.2) {$x-L$};
    \node[below] at (2,-0.2) {$x+L$};
    
    % Arrows
    \draw[->, thick, blue] (-1.8, 0.2) to[bend left] node[midway, above] {$1/2$} (-0.2, 0.2);
    \draw[->, thick, blue] (1.8, 0.2) to[bend right] node[midway, above] {$1/2$} (0.2, 0.2);
    
    \node at (0, 1.5) {\textbf{משוואת המאסטר (Master Equation):}};
    \node at (0, 1) {$P(x, t+\Delta t) = \frac{1}{2}P(x-L, t) + \frac{1}{2}P(x+L, t)$};
\end{tikzpicture}
\end{center}

כדי להגיע למשוואה דיפרנציאלית, נבצע \textbf{גבול הרצף} ($L \to 0, \Delta t \to 0$) באמצעות פיתוח טור טיילור.

\textbf{אגף שמאל (פיתוח בזמן):}
\[
P(x, t+\Delta t) \approx P(x,t) + \Delta t \frac{\partial P}{\partial t}
\]

\textbf{אגף ימין (פיתוח במרחב):}
נפתח את האיברים $P(x\pm L, t)$ עד סדר שני ב-$L$ (סדר ראשון יתבטל בגלל הסימטריה):
\[
P(x\pm L, t) \approx P(x,t) \pm L \frac{\partial P}{\partial x} + \frac{L^2}{2} \frac{\partial^2 P}{\partial x^2}
\]
נציב במשוואת המאסטר:
\begin{align*}
P + \Delta t \frac{\partial P}{\partial t} &= \frac{1}{2}\left( P - L P_x + \frac{L^2}{2} P_{xx} \right) + \frac{1}{2}\left( P + L P_x + \frac{L^2}{2} P_{xx} \right) \\
P + \Delta t \frac{\partial P}{\partial t} &= P + 0 + \frac{L^2}{2} \frac{\partial^2 P}{\partial x^2}
\end{align*}
נצמצם את $P(x,t)$ משני הצדדים ונחלק ב-$\Delta t$:
\[
\frac{\partial P}{\partial t} = \left( \frac{L^2}{2\Delta t} \right) \frac{\partial^2 P}{\partial x^2}
\]

כדי שהגבול יהיה מוגדר היטב, עלינו לדרוש שהיחס $\frac{L^2}{2\Delta t}$ ישאף לקבוע סופי כאשר $L,\Delta t \to 0$.

\begin{resultbox}{משוואת הדיפוזיה מתוך הלוך אקראי}
נגדיר את מקדם הדיפוזיה $D = \lim_{L,\Delta t \to 0} \frac{L^2}{2\Delta t}$. נקבל:
\[
\frac{\partial P}{\partial t} = D \frac{\partial^2 P}{\partial x^2}
\]
\end{resultbox}

\subsection{תכונות בסיסיות של משוואת הדיפוזיה}

\begin{theorembox}{החלקה (Smoothing)}
בניגוד למשוואת הגלים שמשמרת את צורת הגל, משוואת הדיפוזיה "ממיסה" את הפרופיל ההתחלתי.
\begin{itemize}
    \item פיקים חדים דועכים ומתרחבים.
    \item בזמן $t \to \infty$, הפתרון שואף למצב אחיד (קבוע) או לאפס (תלוי בתנאי השפה).
    \item אינפורמציה על המצב ההתחלתי המדויק אובדת עם הזמן (תהליך בלתי הפיך, הקשור לעלייה באנטרופיה).
\end{itemize}
\end{theorembox}

\subsubsection{חוקי שימור תחת דיפוזיה}
נסתכל על המסה הכוללת (או האנרגיה הכוללת) במערכת אינסופית חד-ממדית: $M(t) = \int_{-\infty}^{\infty} u(x,t) \, dx$.
נבדוק את הנגזרת בזמן:
\[
\frac{dM}{dt} = \int_{-\infty}^{\infty} \frac{\partial u}{\partial t} \, dx = \int_{-\infty}^{\infty} D \frac{\partial^2 u}{\partial x^2} \, dx = D \left[ \frac{\partial u}{\partial x} \right]_{-\infty}^{\infty}
\]
עבור מערכת אינסופית שבה הפונקציה והנגזרת שלה דועכות לאפס בקצוות ($\lim_{x\to\pm\infty} u_x = 0$), נקבל $\frac{dM}{dt} = 0$.
\textbf{מסקנה:} המסה הכוללת נשמרת.

באופן דומה, ניתן להראות ש\textbf{מרכז המסה} $\langle x \rangle = \frac{1}{M} \int x u \, dx$ נשמר בזמן (לא זז) עבור דיפוזיה סימטרית ללא סחיפה (\LR{Drift}). זה נובע מכך שלכל חלקיק שזז ימינה יש בממוצע חלקיק שזז שמאלה.

